\documentclass[11pt]{article}

\usepackage[margin=1in]{geometry}
\usepackage{amsmath, amssymb, amsthm, bm, mathtools}
\usepackage{enumitem}
\usepackage{hyperref}

\newcommand{\bY}{\mathbf{Y}}
\newcommand{\bX}{\mathbf{X}}
\newcommand{\bD}{\mathbf{D}}
\newcommand{\bF}{\mathbf{F}}
\newcommand{\bH}{\mathbf{H}}
\newcommand{\bG}{\mathbf{G}}
\newcommand{\bM}{\mathbf{M}}
\newcommand{\bR}{\mathbf{R}}
\newcommand{\bA}{\mathbf{A}}
\newcommand{\bS}{\mathbf{S}}
\newcommand{\bP}{\mathbf{P}}
\newcommand{\bz}{\mathbf{z}}
\newcommand{\br}{\mathbf{r}}
\newcommand{\bzero}{\mathbf{0}}
\newcommand{\diag}{\operatorname{diag}}
\newcommand{\tr}{\operatorname{tr}}
\newcommand{\cov}{\operatorname{cov}}
\newcommand{\var}{\operatorname{var}}
\newcommand{\E}{\operatorname{E}}
\newcommand{\I}{\mathbf{I}}

\title{A Proposed MRTS-Augmented Bayesian Hierarchical Skew-$t$ Model}
\author{}
\date{\today}

\begin{document}

\maketitle

\section{Goal}

The goal of this note is to formulate a computationally efficient alternative to
the spatial dependence component in the Morris skew-$t$ hierarchical model. The
main idea is to retain the original hierarchical structure for the mean, skew,
and scale components, but replace the site-level Mat\'ern correlation model with
an MRTS-based low-rank spatial random effect representation. To avoid treating
the full $K \times K$ covariance matrix $\bM$ as a high-dimensional Bayesian
parameter block, we propose to update $\bM$ using the closed-form maximum
likelihood estimators derived by Tzeng and Huang (2018), and then convert the
resulting covariance matrix to a correlation matrix so that the original Morris
separation between marginal scale and spatial dependence is preserved.

\section{Original Morris Hierarchical Structure}

Let $s_1, \ldots, s_n \in D \subset \mathbb{R}^d$ denote the observed spatial
locations and let $t = 1, \ldots, T$ index replicates in time. Let
$Y_t(s_i)$ denote the response at location $s_i$ and replicate $t$. For each
$t$, define the response vector
\[
\bY_t = \bigl(Y_t(s_1), \ldots, Y_t(s_n)\bigr)'.
\]

As in Morris et al., we retain the partition-based skew and scale structure. Let
$P_{tk}$, $k = 1, \ldots, K_t^{(p)}$, denote the spatial partitions at time $t$.
For $s \in P_{tk}$, define
\[
z_t(s) = z_{tk},
\qquad
\sigma_t^2(s) = \sigma_{tk}^2.
\]
Let
\[
\bz_t = \bigl(z_t(s_1), \ldots, z_t(s_n)\bigr)',
\qquad
\bD_t = \diag\bigl(\sigma_t(s_1), \ldots, \sigma_t(s_n)\bigr).
\]
Let $\bX_t$ denote the $n \times p$ design matrix for replicate $t$, and let
$\bm{\beta} \in \mathbb{R}^p$ and $\lambda \in \mathbb{R}$ denote the regression
and skewness parameters.

The original Morris model can be written in vector form as
\begin{equation}
\bY_t = \bX_t \bm{\beta} + \lambda \bD_t |\bz_t| + \bD_t \mathbf{v}_t,
\qquad
t = 1, \ldots, T,
\label{eq:morris}
\end{equation}
where $|\bz_t|$ is interpreted componentwise, and $\mathbf{v}_t$ is a
zero-mean Gaussian spatial process with a correlation matrix determined by a
Mat\'ern family with a nugget-type parameter. In the original implementation,
this correlation structure is governed by a small number of hyperparameters
such as $\rho$, $\nu$, and $\gamma$, but repeated MCMC updates still require
repeated inversions of an $n \times n$ covariance matrix.

\section{Proposed MRTS Spatial Layer}

\subsection{Basis Expansion}

Let $f_1(\cdot), \ldots, f_K(\cdot)$ denote prespecified MRTS basis functions,
with $K \ll n$ in the intended applications. Define
\[
\mathbf{f}(s) = \bigl(f_1(s), \ldots, f_K(s)\bigr)',
\]
and let the $n \times K$ basis matrix be
\[
\bF = \begin{bmatrix}
\mathbf{f}(s_1)' \\
\vdots \\
\mathbf{f}(s_n)'
\end{bmatrix}.
\]

Introduce a latent spatial random effect vector
\[
\mathbf{u}_t = \bF \bm{\eta}_t + \bm{\xi}_t,
\qquad
\bm{\eta}_t \sim N_K(\bzero, \bM),
\qquad
\bm{\xi}_t \sim N_n(\bzero, \sigma_\xi^2 \I_n),
\]
independently over $t$, where $\bM$ is an unknown nonnegative-definite
$K \times K$ matrix and $\sigma_\xi^2 \ge 0$ is a fine-scale variance
parameter. Conditional on $(\bM, \sigma_\xi^2)$,
\[
\mathbf{u}_t \sim N_n(\bzero, \bG),
\qquad
\bG = \bF \bM \bF' + \sigma_\xi^2 \I_n.
\]

\subsection{Standardization to a Correlation Matrix}

If $\bG$ were used directly in \eqref{eq:morris}, then the marginal variance of
the spatial term would be controlled jointly by $\bG$ and by $\bD_t$, which
would blur the original division of labor between the scale process and the
spatial dependence process. To preserve the Morris interpretation, we convert
$\bG$ to a correlation matrix.

Define
\[
\bA = \diag(\bG),
\qquad
\bR = \bA^{-1/2} \bG \bA^{-1/2}.
\]
Then $\bR$ is positive semidefinite with unit diagonal. We define the
standardized latent spatial effect
\[
\mathbf{v}_t = \bA^{-1/2} \mathbf{u}_t,
\qquad
\mathbf{v}_t \sim N_n(\bzero, \bR).
\]
The proposed hierarchical model is therefore
\begin{equation}
\bY_t = \bX_t \bm{\beta} + \lambda \bD_t |\bz_t| + \bD_t \mathbf{v}_t,
\qquad
\mathbf{v}_t \sim N_n(\bzero, \bR),
\label{eq:proposed}
\end{equation}
with
\[
\bR = \bA^{-1/2} (\bF \bM \bF' + \sigma_\xi^2 \I_n) \bA^{-1/2},
\qquad
\bA = \diag(\bF \bM \bF' + \sigma_\xi^2 \I_n).
\]

\subsection{Conditional Residual Representation}

Given $(\bm{\beta}, \lambda, \bz_t, \bD_t)$, define the standardized residual
vector
\begin{equation}
\br_t
=
\bD_t^{-1} \bigl( \bY_t - \bX_t \bm{\beta} - \lambda \bD_t |\bz_t| \bigr),
\qquad
t = 1, \ldots, T.
\label{eq:resid}
\end{equation}
Under the proposed model,
\[
\br_t = \mathbf{v}_t \sim N_n(\bzero, \bR).
\]

For the purpose of updating the low-rank covariance structure, we also consider
the unstandardized low-rank covariance
\[
\bG = \bF \bM \bF' + \sigma_\xi^2 \I_n,
\]
and use the Tzeng--Huang closed-form estimators as a plug-in update for
$(\bM, \sigma_\xi^2)$ before converting $\bG$ to $\bR$.

\section{Closed-Form Update for \texorpdfstring{$\bM$ and $\sigma_\xi^2$}{M and sigma-xi-squared}}

\subsection{Working Covariance Update}

At a given MCMC iteration, treat the residual vectors $\br_1, \ldots, \br_T$ in
\eqref{eq:resid} as the current approximately Gaussian replicates used to update
the spatial dependence block. Define the empirical residual covariance matrix
\[
\bS = \frac{1}{T} \sum_{t=1}^T \br_t \br_t'.
\]

The spatial random effect model of Tzeng and Huang takes the form
\[
\mathbf{z}_t = \mathbf{y}_t + \bm{\epsilon}_t,
\qquad
\mathbf{y}_t = \bF \bm{\eta}_t + \bm{\xi}_t,
\]
with
\[
\bm{\eta}_t \sim N_K(\bzero, \bM),
\qquad
\bm{\xi}_t \sim N_n(\bzero, \sigma_\xi^2 \I_n),
\qquad
\bm{\epsilon}_t \sim N_n(\bzero, \sigma_\epsilon^2 \I_n).
\]
In the present proposal, there is no separate measurement-error term in the
standardized residual update, so the natural choice is
\[
\sigma_\epsilon^2 = 0.
\]
If a known measurement-error variance is desired in a specific application, it
can be inserted without changing the formulas below.

\subsection{Theorem 2 Plug-In Estimator}

Define
\[
\bH = (\bF' \bF)^{-1/2} \bF' \bS \bF (\bF' \bF)^{-1/2}.
\]
Let the eigendecomposition of $\bH$ be
\[
\bH = \bP \diag(d_1, \ldots, d_K) \bP',
\]
where $d_1 \ge \cdots \ge d_K \ge 0$ are the eigenvalues. Let $d_0 = 0$. Define
$L^*$ as the largest $L \in \{1, \ldots, K\}$ such that
\[
d_L >
\max \left\{
\frac{1}{n-L}
\left(
\tr(\bS) - \sum_{k=1}^L d_k
\right),
\sigma_\epsilon^2
\right\},
\]
if such an $L$ exists, and let $L^* = 0$ otherwise.

Then the Tzeng--Huang estimators are
\begin{equation}
\widehat{\sigma}_{\xi}^2
=
\max \left\{
\frac{1}{n-L^*}
\left(
\tr(\bS) - \sum_{k=1}^{L^*} d_k
\right)
- \sigma_\epsilon^2,
\,
0
\right\},
\label{eq:sigxi}
\end{equation}
and
\begin{equation}
\widehat{\bM}
=
(\bF' \bF)^{-1/2}
\bP
\diag(\hat d_1, \ldots, \hat d_K)
\bP'
(\bF' \bF)^{-1/2},
\label{eq:mhat}
\end{equation}
where
\[
\hat d_k = \max(d_k - \widehat{\sigma}_{\xi}^2 - \sigma_\epsilon^2, 0),
\qquad
k = 1, \ldots, K.
\]

The updated low-rank covariance is then
\begin{equation}
\widehat{\bG} = \bF \widehat{\bM} \bF' + \widehat{\sigma}_{\xi}^2 \I_n.
\label{eq:ghat}
\end{equation}
To preserve the Morris architecture, define
\begin{equation}
\widehat{\bA} = \diag(\widehat{\bG}),
\qquad
\widehat{\bR} = \widehat{\bA}^{-1/2} \widehat{\bG} \widehat{\bA}^{-1/2}.
\label{eq:rhat}
\end{equation}

\subsection{Interpretation}

The update in \eqref{eq:sigxi}--\eqref{eq:rhat} does \emph{not} treat the
$K(K+1)/2$ free elements of $\bM$ as a Bayesian parameter block requiring its
own MCMC updates. Instead, it profiles the covariance block out by replacing
$(\bM, \sigma_\xi^2)$ with closed-form estimators computed from the current
residual replicates. Thus the computational role of the MRTS layer is to reduce
the dimension of the matrix decomposition from $n \times n$ to $K \times K$,
rather than to introduce a large new stochastic parameter block.

\section{Likelihood and Computation}

\subsection{Conditional Gaussian Likelihood}

Given $(\bm{\beta}, \lambda, \bz_t, \bD_t, \bR)$, the conditional distribution
of $\bY_t$ is
\[
\bY_t \mid \bm{\beta}, \lambda, \bz_t, \bD_t, \bR
\sim
N_n\bigl(
\bX_t \bm{\beta} + \lambda \bD_t |\bz_t|,
\,
\bD_t \bR \bD_t
\bigr).
\]
Equivalently, the standardized residuals satisfy
\[
\br_t \mid \bR \sim N_n(\bzero, \bR),
\qquad
t = 1, \ldots, T.
\]
Hence, up to an additive constant, the conditional log-likelihood of the
residual block is
\begin{equation}
\ell(\bR \mid \br_1, \ldots, \br_T)
=
-\frac{T}{2} \log |\bR|
-\frac{1}{2} \sum_{t=1}^T \br_t' \bR^{-1} \br_t.
\label{eq:loglik}
\end{equation}

\subsection{Woodbury Identity for Fast Matrix Operations}

Let
\[
\bG = \bF \bM \bF' + \sigma_\xi^2 \I_n,
\qquad
\sigma_\xi^2 > 0.
\]
Then the Woodbury identity gives
\begin{equation}
\bG^{-1}
=
\sigma_\xi^{-2} \I_n
- \sigma_\xi^{-2} \bF
\left(
\bM^{-1} + \sigma_\xi^{-2} \bF' \bF
\right)^{-1}
\bF' \sigma_\xi^{-2}.
\label{eq:woodbury}
\end{equation}
Also, by the matrix determinant lemma,
\begin{equation}
\log |\bG|
=
n \log(\sigma_\xi^2)
+ \log |\bM|
+ \log \left|
\bM^{-1} + \sigma_\xi^{-2} \bF' \bF
\right|.
\label{eq:detg}
\end{equation}

Because
\[
\bR = \bA^{-1/2} \bG \bA^{-1/2},
\]
we obtain
\begin{equation}
\bR^{-1} = \bA^{1/2} \bG^{-1} \bA^{1/2},
\label{eq:rinv}
\end{equation}
and
\begin{equation}
\log |\bR| = \log |\bG| - \sum_{i=1}^n \log(A_{ii}).
\label{eq:logdetr}
\end{equation}
Thus all large-matrix computations needed in \eqref{eq:loglik} can be reduced to
operations involving $\bF' \bF$ and other $K \times K$ matrices, followed by
matrix-vector products with $\bF$.

\subsection{Computational Complexity}

For a single covariance update, the main costs are:
\begin{itemize}[leftmargin=2em]
\item forming $\bF' \br_t$ for $t = 1, \ldots, T$, which costs
$O(nKT)$;
\item forming the $K \times K$ matrix
$\bF' \bS \bF = T^{-1} \sum_{t=1}^T (\bF' \br_t)(\bF' \br_t)'$, which costs
$O(K^2 T)$ after the projected residuals are available;
\item eigendecomposition or inversion of $K \times K$ matrices, which costs
$O(K^3)$.
\end{itemize}
This replaces repeated direct inversions of dense $n \times n$ Mat\'ern
covariance matrices, whose dominant cost is typically $O(n^3)$ per update.
When $K \ll n$, this is the main source of computational acceleration.

\subsection{Identifiability Constraint Imposed by \texorpdfstring{$T$}{T}}

The computational gain from choosing $K \ll n$ does not remove the statistical
constraint imposed by the number of replicates $T$. To see this, define the
projected residual vectors
\[
\widetilde{\br}_t = (\bF' \bF)^{-1/2} \bF' \br_t \in \mathbb{R}^K,
\qquad
t = 1, \ldots, T.
\]
Then the matrix whose eigenvalues drive the Tzeng--Huang update is
\begin{equation}
\bH
=
(\bF' \bF)^{-1/2} \bF' \bS \bF (\bF' \bF)^{-1/2}
=
\frac{1}{T}
\sum_{t=1}^T
\widetilde{\br}_t \widetilde{\br}_t'.
\label{eq:Hdecomp}
\end{equation}
Therefore,
\begin{equation}
\operatorname{rank}(\bH) \le \min(K, T).
\label{eq:rankbound}
\end{equation}
If centered residuals are used in place of raw residual replicates, the bound is
typically reduced to $\operatorname{rank}(\bH) \le \min(K, T-1)$.

Equation \eqref{eq:rankbound} implies that even when the replicates are
conditionally independent and identically distributed under a common covariance
matrix, the effective rank that can be learned for $\bM$ is still limited by the
sample size in time. In particular:
\begin{itemize}[leftmargin=2em]
\item if $K > T$, then at most $T$ directions in the $K$-dimensional basis space
can be identified from the replicate covariance;
\item if $K$ is close to $T$, then the leading eigendirections may still be
estimable, but the smaller eigenvalues and associated eigenspaces are expected
to be unstable;
\item the theorem still provides a valid closed-form update, but the resulting
$\widehat{\bM}$ will be effectively low-rank and may be noisy in finite samples.
\end{itemize}

Consequently, basis selection should be guided not only by the spatial sample
size $n$, but also by the number of temporal replicates $T$. A practical
recommendation for the first implementation is to choose $K$ so that $K \ll T$
in addition to $K \ll n$, and to regard the Tzeng--Huang truncation step as an
adaptive effective-rank selector rather than as a license to let $K$ grow
without regard to $T$.

\section{Hybrid Bayesian--Profile Inference Scheme}

The proposed method is not a fully Bayesian model for $\bM$. Instead, it is a
Bayesian hierarchical model for the skew, scale, regression, and partition
blocks, combined with a profiled low-rank covariance update. A practical
algorithm is:

\begin{enumerate}[leftmargin=2em]
\item Update latent censored or missing values in $\bY_t$ as in the original
Morris computation.
\item Update $(\bm{\beta}, \lambda)$ using their conditional Gaussian blocks.
\item Update the partition-specific latent variables $\bz_t$, the scale terms
$\bD_t$, and any partition-location parameters exactly as in the original Morris
sampler.
\item Compute residuals $\br_t$ using \eqref{eq:resid}.
\item Form $\bS = T^{-1} \sum_{t=1}^T \br_t \br_t'$.
\item Update $(\bM, \sigma_\xi^2)$ by the closed-form formulas
\eqref{eq:sigxi}--\eqref{eq:mhat}.
\item Construct $\bG$, $\bA$, and $\bR$ via \eqref{eq:ghat}--\eqref{eq:rhat}.
\item Evaluate quadratic forms and log-determinants using
\eqref{eq:woodbury}--\eqref{eq:logdetr}.
\end{enumerate}

This scheme is best viewed as a hybrid MCMC/ECM algorithm: the Morris blocks are
updated by Bayesian conditional simulation, whereas the low-rank covariance
block is updated by a closed-form optimization step conditional on the current
residuals.

\subsection{Implementation Checklist}

Before coding the first working version, the following implementation choices
should be fixed:
\begin{enumerate}[leftmargin=2em]
\item \textbf{Fixed basis design.} The MRTS basis matrix $\bF$ is treated as
fixed throughout the analysis and is constructed outside the MCMC.
\item \textbf{Time-invariant covariance.} The residual replicates
$\br_1, \ldots, \br_T$ are assumed conditionally independent and identically
distributed under a common spatial dependence matrix.
\item \textbf{Profiled covariance block.} The pair $(\bM, \sigma_\xi^2)$ is
updated by the closed-form Tzeng--Huang step rather than by a high-dimensional
Bayesian parameter update.
\item \textbf{Correlation normalization.} The covariance $\bG$ is converted to a
correlation matrix $\bR$ before entering the Morris likelihood, so that
$\bD_t$ remains responsible for marginal scale.
\item \textbf{Numerical stabilization.} The implementation should include a
small ridge or floor when computing $(\bF' \bF)^{-1/2}$ and when enforcing
$\widehat{\sigma}_\xi^2 \ge 0$.
\item \textbf{Basis size relative to $T$.} The chosen basis dimension $K$ should
respect the identifiability constraint described in
\eqref{eq:Hdecomp}--\eqref{eq:rankbound}.
\item \textbf{Validation targets.} The first implementation should be evaluated
against the original Morris model using both runtime and predictive criteria,
including uncertainty calibration and tail behavior.
\end{enumerate}

\subsection{Suggested Module Decomposition}

A clean first implementation can be organized around the following modules:
\begin{enumerate}[leftmargin=2em]
\item \texttt{build\_basis\_matrix(...)}:
construct the fixed MRTS basis matrix $\bF$ at the observed locations, along
with any stored quantities such as $\bF' \bF$.
\item \texttt{update\_residuals(...)}:
compute the conditional standardized residuals
$\br_t = \bD_t^{-1}(\bY_t - \bX_t \bm{\beta} - \lambda \bD_t |\bz_t|)$.
\item \texttt{update\_M\_closed\_form(...)}:
compute $\bS$, form $\bH$, carry out the eigendecomposition, determine $L^*$,
and return $(\widehat{\bM}, \widehat{\sigma}_\xi^2)$.
\item \texttt{build\_R\_from\_G(...)}:
form $\widehat{\bG} = \bF \widehat{\bM} \bF' + \widehat{\sigma}_\xi^2 \I_n$,
compute $\widehat{\bA} = \diag(\widehat{\bG})$, and return
$\widehat{\bR} = \widehat{\bA}^{-1/2} \widehat{\bG} \widehat{\bA}^{-1/2}$.
\item \texttt{quadform\_logdet(...)}:
evaluate the quadratic forms and log-determinants needed in the residual
likelihood using the Woodbury identity and matrix determinant lemma.
\item \texttt{update\_morris\_blocks(...)}:
retain the existing Morris updates for censored data imputation, regression,
skewness, partition-specific scale terms, and any knot-location parameters.
\item \texttt{predict\_at\_new\_sites(...)}:
construct basis values at prediction sites and produce posterior predictive
draws using the fitted low-rank spatial layer and the current Morris mean/scale
components.
\end{enumerate}

This modular decomposition is intended to make the first implementation easy to
debug: the low-rank covariance block can be tested first in a simplified
Gaussian setting before it is reattached to the full censored skew-$t$ sampler.

\section{Optional One-Parameter Bayesian Extension}

If some uncertainty in the covariance block is desired without returning to a
full $K(K+1)/2$-dimensional MCMC update, one simple extension is to introduce a
scalar expansion parameter $c > 0$ and set
\[
\bM = c \, \widehat{\bM},
\]
where $\widehat{\bM}$ is the Tzeng--Huang plug-in estimator. One may then place
a prior on $c$, for example
\[
c \sim \operatorname{Gamma}(a_c, b_c),
\]
and update only $c$ within the MCMC. This preserves the dominant eigenstructure
estimated from the closed-form step while allowing limited posterior
uncertainty in the overall strength of the low-rank dependence.

\section{Modeling Assumptions and Practical Remarks}

\begin{enumerate}[leftmargin=2em]
\item The MRTS basis matrix $\bF$ is treated as fixed and prespecified.
\item The Tzeng--Huang update relies on replicate residual vectors
$\br_1, \ldots, \br_T$ that are conditionally independent with a common spatial
dependence structure.
\item In the present proposal, $\sigma_\epsilon^2$ is set to zero because the
standardized residual vectors are treated as directly observed for the purpose of
the covariance update. If an external measurement-error variance is known, it can
be inserted into the update formulas.
\item Standardizing $\bG$ to $\bR$ is a deliberate modeling choice made to
preserve the original Morris interpretation that $\bD_t$ controls marginal scale
while the spatial layer controls dependence.
\item The basis dimension $K$ should remain well below $n$, but it must also be
chosen with care relative to $T$ because the effective rank learnable from the
replicate covariance is constrained by \eqref{eq:rankbound}.
\end{enumerate}

\section{Summary}

The proposed model keeps the Morris skew-$t$ hierarchy intact at the observation
level,
\[
\bY_t = \bX_t \bm{\beta} + \lambda \bD_t |\bz_t| + \bD_t \mathbf{v}_t,
\]
but replaces the site-level Mat\'ern dependence with an MRTS low-rank spatial
layer. The unknown covariance matrix $\bM$ is not sampled as a large Bayesian
parameter block. Instead, it is updated in closed form from current residuals
using the Tzeng--Huang estimator, after which the implied covariance matrix is
standardized to a correlation matrix and inserted back into the Morris
hierarchy. This retains computational scalability while preserving the original
interpretation of the skew and scale blocks.

\end{document}
